
@@M0@@
\uline{est
exacte}.

Si @@M1@@ est un @@M2@@-Module cohérent, @@M3@@ est un @@M4@@-Module dont le support est contenu dans @@M5@@; par conséquent sa restriction à @@M6@@ est nulle, et @@M7@@, d'où, les éléments de la forme @@M8@@ engendrant @@M9@@, @@M10@@. On en déduit un homomorphisme
@@M11@@
On va définir un homomorphisme inverse de @@M12@@. Soit @@M13@@ un @@M14@@-Module cohérent. Comme @@M15@@ est noethérien, il existe un @@M16@@-Module cohérent @@M17@@ dont la restriction à @@M18@@ est isomorphe à @@M19@@. Si @@M20@@ est un second @@M21@@-Module cohérent dont la restriction à @@M22@@ est isomorphe à @@M23@@, alors on a dans @@M24@@ :
@@M25@@
Pour le voir, on considère @@M26@@ comme un sous-faisceau de @@M27@@ au moyen de l'application diagonale; @@M28@@ se prolonge alors en un sous-faisceau @@M29@@ de @@M30@@, et on est ramené à montrer que
@@M31@@
et l'analogue pour @@M32@@. Soient alors @@M33@@ et @@M34@@ le noyau et le conoyau de l'homomorphisme composé :
@@M35@@
cet homomorphisme se réduit à l'identité de @@M36@@ sur @@M37@@, et par suite @@M38@@ et @@M39@@ sont à support contenu dans @@M40@@. Si @@M41@@ est l'idéal de @@M42@@ dans @@M43@@, il en résulte qu'il existe des entiers @@M44@@ et @@M45@@ tels que @@M46@@, @@M47@@ (EGA I 9.3.4). On peut donc dévisser @@M48@@ et @@M49@@ en extensions de Modules
